%=====================================================================
% FRONT MATTER
%=====================================================================
\pagenumbering{roman}
% Front-matter figures sit outside any chapter, so give them their own prefix
% rather than letting them inherit a misleading chapter number.
\structuralfigures{F.}
\thispagestyle{empty}

\begin{center}
\vspace*{1.2cm}

{\color{secondary}\rule{0.85\textwidth}{2pt}}\\[6pt]
{\large\textsc{\color{secondary}Bachelor of Science in Information Technology}}\\[18pt]

{\Huge\bfseries\color{primary}DATA SCIENCE}\\[10pt]
{\LARGE\bfseries\color{primary}From Foundations to Intelligent Systems}\\[16pt]

{\color{secondary}\rule{0.85\textwidth}{2pt}}\\[16pt]

{\Large\textsc{Unified Lecture Handout}}\\[6pt]
{\large\color{gray!70!black}Complete Course in One Volume $\bullet$ Fall 2026}\\[24pt]

% --- The four application domains, announced on the title page ---
\begin{tikzpicture}[
  dom/.style={rounded corners=4pt, draw=#1, line width=1.1pt, fill=#1!6,
              text=#1, align=center, font=\small\bfseries,
              minimum width=3.35cm, minimum height=1.5cm}]
  \node[dom=lensLA]  at (-5.1,0) {\faIcon{graduation-cap}\\[2pt]Learning\\Analytics};
  \node[dom=lensPM]  at (-1.7,0) {\faIcon{tasks}\\[2pt]Project\\Management};
  \node[dom=lensIOT] at ( 1.7,0) {\faIcon{microchip}\\[2pt]Internet\\of Things};
  \node[dom=lensSEC] at ( 5.1,0) {\faIcon{shield-alt}\\[2pt]Cyber\\security};
\end{tikzpicture}\\[10pt]
{\small\itshape\color{gray!70!black}Every chapter shows its methods at work in all four domains.}

\vspace{1.4cm}

\begin{tcolorbox}[enhanced, width=0.86\textwidth, colback=lightbg,
  colframe=primary, boxrule=0pt, leftrule=4pt, arc=2pt, top=6pt, bottom=6pt]
\small
\textbf{\color{primary}What this volume replaces.} Previous offerings of this course
distributed thirteen separate handouts. This edition consolidates them into a
single, resequenced volume in which every topic appears exactly once, in the
order you need it, with a complete set of figures.
\end{tcolorbox}

\vspace{1.2cm}
{\small\color{gray!60!black}Six Parts $\bullet$ Thirty-Two Chapters $\bullet$ Seven Appendices}
\vspace{1cm}
\end{center}

\newpage

%=====================================================================
\chapter*{How to Use This Handout}
\addcontentsline{toc}{chapter}{How to Use This Handout}
\markboth{How to Use This Handout}{}

This is not a reference manual to be searched; it is a \term{path} to be walked.
The chapters are ordered so that every idea rests on ideas you have already met.
Reading out of order is possible, but each chapter tells you honestly what it
assumes.

\section*{The shape of the book}

The course moves through six parts, and the movement is deliberate:

\begin{tight}
  \item \textbf{Part I --- Foundations of Data.} What data science is, what data
        \emph{is}, and the mathematics and statistics you need to describe it.
  \item \textbf{Part II --- From Data to Insight.} Cleaning, exploring,
        visualising, and reasoning honestly from a sample to a population.
  \item \textbf{Part III --- Machine Learning Core.} Learning from data:
        regression, classification, evaluation, ensembles, and finding structure
        without labels.
  \item \textbf{Part IV --- Deep Learning and Modern AI.} Neural networks, deep
        architectures, sequences, language, generative models, and agents.
  \item \textbf{Part V --- Data Science in Systems Context.} Engineering data at
        scale, cloud and edge, and the two application chapters this course is
        built towards: \emph{Data Science for IoT} and \emph{Data Science for
        Cybersecurity}.
  \item \textbf{Part VI --- Professional Practice.} Ethics, research method,
        project management, communication, and a capstone that ties everything
        together.
\end{tight}

\begin{conceptbox}[Why IoT and Cybersecurity Appear in Part V]
Both are \emph{systems} applications, and both depend on machinery you meet
earlier. IoT analytics needs streaming and time series (\chref{timeseries}) and
edge computing (\chref{cloudedge}). Security analytics needs anomaly detection
(\chref{unsupervised}), evaluation under severe class imbalance
(\chref{evaluation}), and sequence models (\chref{timeseries}). Placing them
earlier would mean hand-waving the methods. Placing them here means you can
actually build them.

You do not, however, wait until Part V to meet these domains. From
\chref{landscape} onwards, every chapter carries a \textbf{Four Lenses} panel
showing that chapter's technique at work in learning analytics, project
management, IoT and cybersecurity. You see the destination from the first page.
\end{conceptbox}

\section*{The furniture in every chapter}

Each chapter is built from the same parts, always in the same order, so you
always know where you are:

\rowcolors{2}{white}{lightbg}
\begin{longtable}{@{}p{4.6cm}p{10.6cm}@{}}
\toprule
\rowcolor{primary!15}
\textbf{Element} & \textbf{What it is for} \\
\midrule
\endhead
\faIcon{bullseye}~\textbf{Learning Outcomes} & What you will be able to \emph{do} after the chapter. Written as actions, not topics. \\
\faIcon{link}~\textbf{Before You Start} & The specific earlier chapters this one leans on. \\
\faIcon{tags}~\textbf{Key Terms} & Vocabulary introduced here. All terms also appear in the glossary (Appendix E). \\
\faIcon{book}~\textbf{Definition} & A precise statement of a term. Learn these exactly. \\
\faIcon{lightbulb}~\textbf{Concept} & The intuition behind a definition --- why it is built this way. \\
\faIcon{calculator}~\textbf{Worked Example} & A complete calculation with real numbers. Work through it with pen and paper. \\
\faIcon{exclamation-circle}~\textbf{Common Pitfalls} & Mistakes that are made every semester. Read these twice. \\
\faIcon{compass}~\textbf{Four Lenses} & The chapter's technique applied in learning analytics, project management, IoT and cybersecurity. \\
\faIcon{flask}~\textbf{Try It Yourself} & A short Python (sometimes R) lab you can run immediately. \\
\faIcon{question-circle}~\textbf{Checkpoint} & Three quick self-tests. Answers in Appendix D. \\
\faIcon{clipboard-check}~\textbf{Chapter Summary} & The chapter compressed to what must be retained. \\
\faIcon{pen-fancy}~\textbf{Review Questions} & Multiple-choice, short-answer and applied questions for assessment practice. \\
\bottomrule
\end{longtable}

\section*{How to study with it}

\begin{tightnum}
  \item \textbf{Read the outcomes first.} They tell you what to look for.
  \item \textbf{Do not skip the worked examples.} Reading a calculation is not
        the same as being able to perform one. Cover the answer and try it.
  \item \textbf{Answer the checkpoint before moving on.} If you cannot, the
        next chapter will be harder than it needs to be.
  \item \textbf{Run the labs.} Fifteen minutes of running code teaches more
        than an hour of reading about it.
  \item \textbf{Use the Four Lenses panels to revise.} Being able to explain a
        method in four different domains is strong evidence you understand it.
\end{tightnum}

\begin{alertbox}[A Note on Notation and Rigour]
This handout uses mathematics where mathematics clarifies --- the update rule
for gradient descent, Bayes' theorem, the definition of precision and recall ---
but it does not prove theorems. Formulas are always accompanied by an
explanation in words and, wherever possible, a numeric example. If a formula
looks intimidating, read the words around it first and return to the symbols
afterwards.
\end{alertbox}

\newpage

%=====================================================================
\chapter*{The Course at a Glance}
\addcontentsline{toc}{chapter}{The Course at a Glance}
\markboth{The Course at a Glance}{}

\dsfigh{%
\resizebox{\textwidth}{!}{%
\begin{tikzpicture}[
  part/.style={rounded corners=3pt, fill=#1, text=white, font=\bfseries\footnotesize,
               minimum width=2.5cm, minimum height=0.72cm, align=center},
  ch/.style={rounded corners=2pt, draw=#1, line width=0.7pt, fill=#1!7, text=#1!60!black,
             font=\scriptsize, minimum width=2.5cm, minimum height=0.52cm, align=center},
  arr/.style={-{Stealth[length=2mm]}, primary!45, line width=0.8pt}
]
\node[part=primary]        (p1) at (0,0)    {Part I\\Foundations};
\node[part=secondary]      (p2) at (3.0,0)  {Part II\\Data to Insight};
\node[part=greenhl]        (p3) at (6.0,0)  {Part III\\ML Core};
\node[part=purplehl]       (p4) at (9.0,0)  {Part IV\\Deep Learning};
\node[part=tealhl]         (p5) at (12.0,0) {Part V\\Systems};
\node[part=goldhl]         (p6) at (15.0,0) {Part VI\\Practice};

\foreach \a/\b in {p1/p2,p2/p3,p3/p4,p4/p5,p5/p6}{\draw[arr] (\a) -- (\b);}

% ---- Chapters under each part
\foreach \i/\t in {1/1. Landscape, 2/2. Lifecycle, 3/3. Data, 4/4. Mathematics,
                   5/5. Descriptive, 6/6. Probability}{
  \pgfmathsetmacro{\yy}{-0.75-(\i-1)*0.62}
  \node[ch=primary] at (0,\yy) {\t};}

\foreach \i/\t in {1/7. Wrangling, 2/8. EDA, 3/9. Inference, 4/10. Causality}{
  \pgfmathsetmacro{\yy}{-0.75-(\i-1)*0.62}
  \node[ch=secondary] at (3.0,\yy) {\t};}

\foreach \i/\t in {1/11. ML Basics, 2/12. Regression, 3/13. Classification,
                   4/14. Evaluation, 5/15. Ensembles, 6/16. Unsupervised}{
  \pgfmathsetmacro{\yy}{-0.75-(\i-1)*0.62}
  \node[ch=greenhl] at (6.0,\yy) {\t};}

\foreach \i/\t in {1/17. Neural Nets, 2/18. Architectures, 3/19. Time Series,
                   4/20. NLP, 5/21. Generative AI, 6/22. Agentic AI}{
  \pgfmathsetmacro{\yy}{-0.75-(\i-1)*0.62}
  \node[ch=purplehl] at (9.0,\yy) {\t};}

\foreach \i/\t in {1/23. Data Engineering, 2/24. Cloud \& Edge, 3/25. IoT,
                   4/26. Cybersecurity, 5/27. MLOps}{
  \pgfmathsetmacro{\yy}{-0.75-(\i-1)*0.62}
  \node[ch=tealhl] at (12.0,\yy) {\t};}
% the two application chapters this course is built towards
\node[rounded corners=2pt, draw=tealhl, line width=1.6pt, fill=tealhl!18,
      text=tealhl!50!black, font=\scriptsize\bfseries, minimum width=2.5cm,
      minimum height=0.52cm] at (12.0,-1.99) {25. IoT};
\node[rounded corners=2pt, draw=tealhl, line width=1.6pt, fill=tealhl!18,
      text=tealhl!50!black, font=\scriptsize\bfseries, minimum width=2.5cm,
      minimum height=0.52cm] at (12.0,-2.61) {26. Cybersecurity};

\foreach \i/\t in {1/28. Ethics, 2/29. Research, 3/30. Projects,
                   4/31. Communicating, 5/32. Capstone}{
  \pgfmathsetmacro{\yy}{-0.75-(\i-1)*0.62}
  \node[ch=goldhl] at (15.0,\yy) {\t};}

% ---- The four lenses running underneath everything
\node[rounded corners=3pt, fill=primary!8, draw=primary!30,
      minimum width=17.3cm, minimum height=0.62cm,
      font=\scriptsize\bfseries, text=primary] at (7.5,-4.62)
      {\faIcon{compass}~~Four Lenses in every chapter:~~
       \textcolor{lensLA}{\faIcon{graduation-cap}~Learning Analytics}~~$\bullet$~~
       \textcolor{lensPM}{\faIcon{tasks}~Project Management}~~$\bullet$~~
       \textcolor{lensIOT}{\faIcon{microchip}~IoT}~~$\bullet$~~
       \textcolor{lensSEC}{\faIcon{shield-alt}~Cybersecurity}};
\end{tikzpicture}}}%
{Course map. The six parts run left to right in the order they are taught; the chapters
in each column are the order within that part. The four application lenses run beneath
the whole course, appearing in every chapter.}{coursemap}

\section*{The Four Running Domains}

Four application domains recur in every chapter of this handout. They were
chosen because between them they cover the situations a BSIT graduate actually
meets, and because each one stresses a different weakness in a data science
method.

\begin{fourlenses}{The Domains Themselves}
\lensLA{\textbf{Learning analytics} applies data science to teaching and
learning: predicting which students are at risk of failing, understanding how
learners move through material, measuring whether an intervention worked.
Its defining difficulty is that the data describes \emph{people}, so fairness,
privacy and the risk of a self-fulfilling prophecy are never far away.}

\lensPM{\textbf{IT project management} applies data science to how work gets
delivered: forecasting whether a sprint will finish, estimating effort,
flagging projects drifting towards failure, allocating people. Its defining
difficulty is small, messy, human-reported data --- and the fact that a
prediction changes the behaviour of the team that hears it.}

\lensIOT{\textbf{The Internet of Things} applies data science to sensor
telemetry: detecting a failing motor from its vibration, forecasting building
energy demand, deciding what to compute on the device and what to send to the
cloud. Its defining difficulties are volume, velocity, missing and drifting
sensors, and a hard limit on the computation available at the edge.}

\lensSEC{\textbf{Cybersecurity} applies data science to defending systems:
spotting an intrusion in network traffic, classifying malware, noticing that
an account is behaving unlike itself. Its defining difficulties are extreme
class imbalance --- attacks are rare --- and a genuine \emph{adversary} who
adapts to whatever model you deploy.}
\end{fourlenses}

\begin{conceptbox}[Why Four Domains and Not One]
A method you have only seen in one setting is a method you have memorised. A
method you can transfer to a new setting is a method you understand. The Four
Lenses panels exist to force that transfer on every single topic --- and the
contrast is instructive in itself. The same anomaly detector that is merely
useful for spotting an unusual exam-submission pattern is safety-critical on a
turbine and adversarially attacked on a firewall.
\end{conceptbox}

\section*{Notation Used in This Handout}

\rowcolors{2}{white}{defbg}
\begin{longtable}{@{}p{2.6cm}p{12.6cm}@{}}
\toprule
\rowcolor{primary!15}
\textbf{Symbol} & \textbf{Meaning} \\
\midrule
\endhead
$n$ & Number of observations (rows, records, examples) in a dataset \\
$p$ or $d$ & Number of features (columns, variables, attributes) \\
$x_i$ & The $i$-th observation; a vector of $p$ feature values \\
$x_{ij}$ & The value of feature $j$ for observation $i$ \\
$y_i$ & The true target (label, outcome) for observation $i$ \\
$\hat{y}_i$ & The model's \emph{prediction} of $y_i$; the hat always means ``estimated'' \\
$X$ & The full feature matrix, $n$ rows by $p$ columns \\
$\bar{x}$ & The sample mean of $x$ \\
$\mu,\ \sigma$ & Population mean and population standard deviation \\
$s$ & Sample standard deviation \\
$P(A)$ & Probability of event $A$ \\
$P(A \mid B)$ & Probability of $A$ \emph{given that} $B$ has occurred \\
$\theta$ or $w$ & Model parameters (weights) to be learned \\
$\alpha$ & Learning rate in gradient descent; also significance level in testing \\
$\lambda$ & Regularisation strength \\
$L(\cdot)$ & Loss function --- how wrong the model is \\
$\nabla$ & Gradient: the vector of partial derivatives \\
$\sum_{i=1}^{n}$ & Sum over all $n$ observations \\
$\|v\|$ & Length (norm) of vector $v$ \\
\bottomrule
\end{longtable}

\begin{alertbox}[Two Symbols Students Confuse Every Year]
$y$ versus $\hat{y}$: $y$ is the \emph{truth}, $\hat{y}$ is the \emph{guess}. Every
error measure in this book is some way of summarising the gaps between them.\\[3pt]
$\sigma$ versus $s$: $\sigma$ describes the whole population (usually unknown);
$s$ is what you compute from your sample as an estimate of it. Confusing them is
the root of most misunderstandings in \chref{inference}.
\end{alertbox}

\section*{Prerequisite Self-Check}

You are ready for Part I if you can answer these. If several give you trouble,
work through \chref{math} slowly and use Appendix C as a companion.

\begin{tightnum}
  \item Compute the average of $4, 8, 8, 10, 15$.
  \item Rewrite $\frac{3}{4}$ as a percentage and as a decimal.
  \item If a bag holds 3 red and 7 blue balls, what is the chance of drawing red?
  \item Solve $2x + 6 = 20$ for $x$.
  \item In a spreadsheet of students, what does one \emph{row} usually represent?
        What does one \emph{column} represent?
  \item Write a line of code, in any language, that prints the numbers 1 to 5.
\end{tightnum}

\smallskip
\textit{Answers: (1)~9\quad (2)~75\%, 0.75\quad (3)~3/10 = 0.3\quad (4)~$x=7$\quad
(5) one row = one student (an observation); one column = one attribute (a feature)\quad
(6) e.g.\ \texttt{for i in range(1,6): print(i)}}

\newpage
\tableofcontents
\newpage
\listoffigures
\addcontentsline{toc}{chapter}{List of Figures}
% No List of Tables: tables in this handout are inline reference material rather
% than numbered exhibits, so they are not captioned and a LoT would be empty.
\newpage

\pagenumbering{arabic}
\setcounter{page}{1}
